\chapter{Developing the method on tank processes}
\label{ch:tanks}
% Merge of the former chapters "Toy examples" (toy_examples.tex) and
% "Setpoint coordination" (coordination.tex), at about half their length
% (THESIS_GUIDE.md, Ch. 4). Purpose: the history that justifies the method,
% and two results that the mill confirms. NOT a second results chapter:
% lessons first, each with its evidence. Target length 8-10 pages.
% The detailed tank material is in Appendix app:tanks (fault-detection
% interface and score, single/two-tank details, full coordination run table)
% and Appendix app:four-tank-protocol (the registered four-tank protocol).
\TODO{Introduction (one paragraph): before the grinding circuit was
available, the method was built on four tank test beds of increasing
difficulty. The first three ask the supervisor to detect leaks (fault
detection, outside the research questions); the fourth rehearses setpoint
coordination against an MPC. Six lessons shaped the method of
Chapter~\ref{ch:method}, and two of them (Sections~\ref{sec:tank-comparator}
and~\ref{sec:tank-robustness}) reappear on the grinding circuit. The tanks
used \texttt{deepseek-flash} and a \emph{held-out} battery that was never used
for promotion.}

\section{The test beds}
\label{sec:tank-testbeds}
% Setup in two sentences each; details in Appendix app:tanks.
\paragraph{Leak detection.} Three processes with leaks as the fault: a
single tank, a two-tank cascade and the quadruple-tank process
\cite{johansson2000} from the PC-Gym benchmark suite \cite{bloor2024pcgym}.
The supervisor raises one anomaly flag per loop and may adjust the setpoints;
the score (Appendix~\ref{app:tanks}) charges tracking error, safety
violations, missed anomalies and false positives.
\TODO{Two sentences on fault detection moved here from the old Background
section: model-based fault detection compares measured behaviour with a model
and alarms on a large residual \cite{isermann2006}; the trade-off between
missed detections, false alarms and delay is reflected in the score.
One-line results: on all three the generated supervisors reduced the
development score by 72--96\% against PID control alone (single tank 76\%,
two-tank 72\%, quadruple tank matched the oracle).}

\paragraph{Setpoint coordination.}

\TODO{Short version of the coordination test bed: the quadruple-tank process
with two cross-paired PI loops; product = production rate $Q$; level $h_2$ in
a band; upper levels below a limit; two setpoints, one target, one degree of
freedom; disturbances on feed, pump gain and valve split as steps, ramps and
sines; 42 development, 42 held-out and 36 beyond-range scenarios; supervisor
every 10~s on a 600~s window; score Equation~\eqref{eq:coord-score}. Keep the
analog table.}
\begin{equation}
  J_s = \mathrm{IAE}_{Q,s} + 2\,(B_s + U_s) + 10\,V_s + 100\,\mathrm{TV}_s
        + 10^3\,E_s ,
  \label{eq:coord-score}
\end{equation}
% IAE_Q: production error (L); B, U: seconds with h2 outside its band / an
% upper level above its limit; V: seconds outside safety bounds; TV: setpoint
% total variation (m); E: exceptions and timeouts.

\begin{table}[tb]
  \centering
  \caption{The grinding-circuit objectives and their analogs on the
    quadruple-tank process.}
  \label{tab:coord-analog}
  \begin{tabular}{@{}ll@{}}
    \toprule
    Grinding circuit & Quadruple-tank analog \\
    \midrule
    P80 on target                  & production rate $Q$ on target \\
    Sump level stable              & $h_2$ inside a band \\
    Circulating load within limits & $h_3$ and $h_4$ below a limit \\
    Fresh feed flow variation      & draw-off at tank 1 or 2, extra inflow at tank 2 \\
    Ore density variation          & loss of pump gain $k_1$ or $k_2$ \\
    Percentage solids variation    & shift of valve split $\gamma_1$ or $\gamma_2$ \\
    Setpoint smoothness            & total variation of the setpoints \\
    \bottomrule
  \end{tabular}
\end{table}

\section{Specification gaming and the regression guard}
\label{sec:tank-gaming}
% Lesson 1 (two-tank cascade).
The two-tank cascade exposed specification gaming: a candidate
improved the average score by flagging anomalies almost constantly. A false
positive cost the same in the fault-free scenario as anywhere else, and the
loss there was outweighed by fewer missed anomalies elsewhere. Other
candidates improved the average while making one scenario much worse; 10 of
58 trials were of this kind. This led to the asymmetric false-positive cost
and to the regression guard (Section~\ref{sec:promotion}).
\TODO{Optionally one or two sentences on what the LLM learned there (the
mitigation that corrupted its own detection statistic, fixed by normalising
pump effort with the level). Otherwise drop.}

\section{Check the test bed before the search}
\label{sec:tank-testbed-check}
% Lesson 2 (quadruple tank, fault detection).
The first training runs on the quadruple-tank process stalled at
$J \approx 36\,300$. The cause was three defects in the test bed rather than
a limit of the search: faults injected through PC-Gym's environment wrapper
were never applied; a decision interval equal to the window length put a
floor on the score; and the diagonal loop pairing ($\lambda = -0.07$) made the
severe tank-1 scenario impossible for any setpoint policy. At the plateau the
champion was within about 7\% of the best score any supervisor could reach on
the flawed test bed. A perfect-detection oracle, which sets each flag to the
true fault state, exposed the plateau; computing such a bound before
interpreting a plateau would have revealed the problem much earlier. After
the corrections, the best supervisors matched the oracle
(Table~\ref{tab:four}).
\TODO{Shorten or expand to taste; the long version is in git history
(\texttt{toy\_examples.tex}).}

\begin{table}[tb]
  \centering
  \caption{Quadruple-tank process (leak detection) after the test-bed
    corrections.}
  \label{tab:four}
  \begin{tabular}{@{}lrr@{}}
    \toprule
    Supervisor & Development & Held-out \\
    \midrule
    PID only (and the seed supervisor)          & 64\,865 & 69\,569 \\
    Plateau champion, not retrained             & 10\,777 & 14\,941 \\
    Best generated (low reasoning, generation 2) &  2\,635 &  2\,609 \\
    Perfect-detection oracle                    &  2\,635 &  2\,609 \\
    \bottomrule
  \end{tabular}
\end{table}

\section{Reasoning is needed; low effort is enough for detection}
\label{sec:tank-reasoning}
% Lesson 3: the only model ablation in the thesis. Keep.
The same search was run from the seed supervisor with three reasoning
settings (Figure~\ref{fig:reasoning} and Table~\ref{tab:reasoning}). With
high and low reasoning the search matched the oracle on both batteries within
three and two generations, respectively, and the held-out scores followed the
development scores. Without reasoning the search improved quickly at first
but stalled at 11\,837 from the third generation on, about 4.5 times the
oracle score. Three weaknesses stood out without reasoning: it missed 49--83
samples in each scenario with a mild leak, against 4--12 with reasoning; its lessons contained
incorrect physics, for example that a leak lowers the other loop's effort
through a shared outflow path that does not exist; and three of twelve
candidates in generations 4--6 called a built-in function that the sandbox
does not provide, after the prompt's list of allowed functions ended in
``etc.''. The list is now generated from the sandbox itself. With high
reasoning, 8 of 15 requests returned an empty response that was still
billed.

\begin{figure}[tb]
  \centering
  \includegraphics[width=\textwidth]{four_tank_reasoning_comparison}
  \caption{Score of the champion after each generation for three reasoning
    settings, on the development and held-out batteries (log scale). Dots
    are individual candidates; the three near $7\cdot 10^5$ crashed on every
    call. Generation~0 is the seed supervisor. The high-reasoning run used
    low reasoning in its third generation and an earlier, longer prompt in
    its first two.}
  \label{fig:reasoning}
\end{figure}

\begin{table}[tb]
  \centering
  \caption{Cost of reaching the oracle with different reasoning settings
    (DeepSeek off-peak prices, US dollars).}
  \label{tab:reasoning}
  \begin{tabular}{@{}lccc@{}}
    \toprule
    Reasoning & Generations & Output tokens & Cost \\
              & to oracle   & per request   & to oracle \\
    \midrule
    High & 3 & $\approx 61\,000$ & $\approx 0.63$ \\
    Low  & 2 & $\approx 23\,000$ & $\approx 0.11$ \\
    None & not reached in 6 & $\approx 2\,000$ & -- \\
    \bottomrule
  \end{tabular}
\end{table}

Each setting was run only once, and the runs differ slightly in their
prompts as noted in Figure~\ref{fig:reasoning}. The results therefore support
the conclusion that reasoning is needed and that low reasoning is enough for
this task, but not a precise ranking of low against high reasoning; this was
not investigated further.

\section{What the search needs}
\label{sec:tank-search}
% Lesson 4: exactly what the mill trainer inherited. One sentence per change;
% the full run table is in Appendix app:tanks (tab:coord-runs).
The coordination runs used the training loop of
Section~\ref{sec:training-loop} with three candidates per generation. Four
parts of that loop were added during these runs, each after a run that
stalled (all runs in Table~\ref{tab:coord-runs}, Appendix~\ref{app:tanks}):
\begin{itemize}
\item \emph{Regression guard by disturbance type, with a schedule.} The
  per-scenario guard of the leak test beds rejected every candidate that beat
  the seed in the first 600~s run (the best scored 258 against 533), so the
  prompt kept showing only the seed.
\item \emph{Failure traces of rejected candidates}, so that the LLM sees why a
  change failed.
\item \emph{One targeted change per candidate.} Without this rule,
  candidates rewrote the supervisor, keeping 16--65\% of the champion's lines,
  and each fixed one scenario while breaking others. With it they kept
  83--99\%.
\item \emph{A measured record instead of unverified lessons.} In one run all
  nine candidates made the same change, which failed, and the lessons saved
  from them recommended it again.
\end{itemize}
\TODO{One or two sentences: the guard schedule gave the first promotion; nine
generations at low effort then promoted nothing; high effort with the
inherited record promoted in each of four generations (effects not
separated). Repeated runs of the same setup ended far apart (development
292--417): a first sign for C3. All coordination runs cost about US\$4.1.}

\section{A comparator is only as good as its tuning}
\label{sec:tank-comparator}
% Lesson 5. Keep tab:coord-main. tab:coord-metrics and tab:coord-kind dropped.
\TODO{Short description of the MPC comparator (written for this thesis:
grid search over setpoint pairs with an internal model of plant and PI loops
and a bias-update disturbance estimate \cite{qin2003}); untuned, tuned with a
random search of 40 configurations on the development battery (a budget
comparable to the LLM's candidates), and with known disturbance as a
reference. Also that the PI loops were retuned ($M_s \le 1.6$) after the runs
and the champions were not retrained.}

Table~\ref{tab:coord-main} compares the best generated supervisor, the
champion of the high-effort run, with the baselines. It reduces the fixed
recipe's development score by 47\%, but it is 67\% worse than the tuned MPC
(265.4 against 158.8). With the original PI loops and against the untuned
MPC, the same champion had come within 1.3\% (206.3 against 203.7). Two
things changed after the runs. First, the MPC was tuned with a budget
comparable to the LLM's, which lowered its development score by 21\%. Second,
the PI loops were retuned. The MPC's internal model includes the loops, so it
adapts to new gains, but the generated supervisor encodes the old regulatory
dynamics as constants (hold times, prediction horizons of its trim, settling
assumptions), and its development score rose by 59 points. Which constants
cause the loss was not investigated further.
\TODO{One sentence on where the deficit is (production error, 230.5 against
119.7~L; recovery about twice as slow) and on beyond range (both degrade by a
factor of about 4.6--4.8, the generated supervisor gives up the constraints to
hold production).}

\begin{table}[tb]
  \centering
  \caption{Average score on the coordination batteries with the retuned PI
    loops (lower is better). The beyond-range battery is never shown to the
    LLM. The generated supervisors were trained with the original PI loops.}
  \label{tab:coord-main}
  \begin{tabular}{@{}lrrr@{}}
    \toprule
    Controller & Development & Held-out & Beyond range \\
    \midrule
    Fixed recipe (seed)                 & 498.3 & 447.0 & 1\,470.5 \\
    Generated, low reasoning effort     & 278.9 & 289.7 & 1\,349.3 \\
    Generated, high reasoning effort    & 265.4 & 274.0 & 1\,223.0 \\
    Untuned MPC                         & 202.0 & 216.6 & 1\,081.5 \\
    Tuned MPC                           & 158.8 & 170.4 &    770.2 \\
    Tuned MPC, known disturbance        & 140.6 & 163.1 &    738.9 \\
    \bottomrule
  \end{tabular}
\end{table}

\section{Nominal ranking does not predict robustness}
\label{sec:tank-robustness}
% Lesson 6: the same finding as on the mill under measurement noise.
% Make tab:coord-stress compact (e.g. drop the stuck-sensor rows).
\TODO{Two or three sentences: conditions none of the controllers was trained
for, on the 84 development and held-out scenarios: level noise of 1 and 3~mm,
plant mismatch, analyser delay, telemetry faults. The tuned MPC more than
doubles its score at 3~mm noise (+156\%) while the generated supervisor
changes by 3\%; its tuning on a noise-free exact model made it sharper and
more fragile. Point forward to Section~\ref{sec:mill-noise-results}.}

\begin{table}[tb]
  \centering
  \caption{Stress tests on the coordination test bed: mean score over the
    development and held-out batteries (84 scenarios), and the change against
    the same controller without stress.}
  \label{tab:coord-stress}
  \small\setlength{\tabcolsep}{3pt}
  \begin{tabular}{@{}lrrrr@{}}
    \toprule
    Condition & Generated & Fixed recipe & Untuned MPC & Tuned MPC \\
    \midrule
    Nominal             & 269.7          & 472.7          & 209.3          & 164.6 \\
    Noise 1~mm          & 263.7 ($-$2\%)   & 472.5 (0\%)    & 198.0 ($-$5\%)   & 241.9 (+47\%) \\
    Noise 3~mm          & 278.1 (+3\%)   & 472.2 (0\%)    & 244.4 (+17\%)  & 421.0 (+156\%) \\
    Mismatch A          & 801.1 (+197\%) & 955.0 (+102\%) & 881.6 (+321\%) & 838.0 (+409\%) \\
    Mismatch B          & 469.4 (+74\%)  & 560.8 (+19\%)  & 506.2 (+142\%) & 432.2 (+163\%) \\
    Analyser delay      & 277.1 (+3\%)   & 472.7 (0\%)    & 209.3 (0\%)    & 164.6 (0\%) \\
    $h_3$ reads NaN     & 439.9 (+63\%)  & 537.9 (+14\%)  & 490.2 (+134\%) & 531.9 (+223\%) \\
    $h_4$ stuck         & 276.5 (+3\%)   & 472.7 (0\%)    & 212.5 (+2\%)   & 173.7 (+6\%) \\
    Production stuck    & 279.9 (+4\%)   & 472.7 (0\%)    & 209.3 (0\%)    & 164.6 (0\%) \\
    \bottomrule
  \end{tabular}
\end{table}

\section{The generated coordination supervisor}
\label{sec:coord-champion}
% Short version; the five-step walkthrough is dropped unless pages allow.
\TODO{One paragraph: a single 261-line function that the LLM turned from the
fixed recipe into a steady-state optimiser with a disturbance estimator
(mass balance, pump-gain estimate), a search over 61 splits of the target
against a cost, and a predictive feedback trim, an MPC-like structure with
every part visible. Its hold time and trim horizon are exactly the constants
that describe the old loops' transient, which is how the loss in
Section~\ref{sec:tank-comparator} could be pointed to in the code.}

\section{The four-tank protocol, registered and not run}
\label{sec:tank-protocol}
\TODO{One paragraph: a protocol (version~2: eight runs of six generations at
high effort against the tuned MPC on a sealed battery of 183 scenarios) was
registered for this rehearsal and never run, because the grinding circuit and
its real comparator became available and the budget went there. Its design
was reused for the grinding-circuit protocol (Section~\ref{sec:protocol}).
The registered text is in Appendix~\ref{app:four-tank-protocol}.}
